\section{Review of Related Literature}

% \subsection{Introduction to Chapter 2}
This chapter examines the literature that most directly informs the present study on hybrid RSSI and CSI-based localization for hidden IoT devices in cluttered indoor environments. Instead of presenting prior work as a disconnected inventory of methods, the review evaluates how recent scholarship addresses four interrelated concerns: non-cooperative device localization, the instability of RSSI under indoor multipath, the precision advantage of CSI-based fingerprinting, and the practical constraints of low-cost edge deployment. This analytical structure is important because the present study does not merely seek another indoor positioning system; it addresses a narrower and more difficult problem in which the target device is hidden, stationary, and potentially non-associative, while the sensing hardware must remain affordable. Accordingly, the discussion emphasizes where prior work is strong, where it remains limited, and how those limitations justify the proposed hybrid adaptive filtering framework.

\subsection{Related Literature}

\subsubsection{Hidden IoT Detection and Non-Cooperative Localization}
Recent scholarship has clarified that hidden-device localization is not simply a variant of generic indoor positioning; rather, it is a security-oriented problem defined by missing cooperation, incomplete prior knowledge, and the need for passive observation. Sharma et al. (2022), through Lumos, significantly advanced this discussion by demonstrating that diverse hidden IoT devices can be identified and localized in unfamiliar environments without assuming a conventional user-tracking setup. This contribution is important because it moves the field away from cooperative localization assumptions and toward adversarially realistic conditions. However, Lumos also reveals a recurring limitation in the literature: systems designed for hidden-device discovery often prioritize detection generality over fine-grained localization using low-cost hardware. In other words, they establish the security relevance of the problem but do not fully solve the engineering problem addressed by the present study.

Ju et al. (2025) strengthen this line of work by showing that channel fingerprints derived from RTS/CTS-triggered RSSI sequences can support the localization of hidden IoT devices in unknown environments. Their study is especially relevant because it validates RF-side signatures when direct device cooperation is unavailable. Yet the broader implication of this work is more important than the result itself: it suggests that the hidden-device problem cannot be solved reliably through network-layer assumptions alone. Even so, both Sharma et al. (2022) and Ju et al. (2025) leave unresolved the question of how such localization can be implemented in multipath-heavy educational interiors using commodity edge hardware. Thus, the literature establishes the problem setting convincingly, but it remains less conclusive on how hidden-device localization can be made simultaneously passive, low-cost, and precise under dense indoor reflections.

\subsubsection{RSSI-Based Indoor Localization Under Multipath Conditions}
RSSI remains attractive in the literature because it is universally available, lightweight to process, and easy to integrate into embedded systems. Nevertheless, recent studies converge on the same critical conclusion: these advantages are offset by weak spatial reliability in cluttered indoor environments. Abdel Ghany et al.\ (2020) compared narrowband LoRa CSI and RSSI in an outdoor IoT measurement campaign and reported stronger robustness for CSI-derived metrics than for RSSI under spatial and temporal variability \cite{ghany2020robustness}. For indoor Wi-Fi, Yang et al.\ (2021) surveyed recent localization scenarios and noted that RSSI-based methods often persist because of deployment convenience rather than peak precision \cite{yang2021survey}, complementing the earlier RSSI--CSI taxonomy of Yang et al.\ (2013) \cite{yang2013rssi2csi}. Leitch et al.\ (2023) reached a comparable conclusion by situating Wi-Fi RSSI among other indoor localization technologies and showing that its practical accessibility does not translate into resilience against indoor interference \cite{leitch2023indoor}.

What makes these studies especially relevant to the present work is not merely their critique of RSSI accuracy, but the implicit design lesson that follows from it. If RSSI is judged only by desk-level precision, it appears inadequate; however, if it is evaluated as a coarse and computationally efficient indicator of spatial proximity, its role becomes strategically useful. This distinction is often underdeveloped in prior studies, which tend to compare RSSI and CSI as competing features rather than as complementary ones. The present study departs from that binary framing. Instead of discarding RSSI because of its multipath sensitivity, it treats RSSI as an efficient first-stage estimator that constrains the search area before finer CSI analysis is invoked. The gap, therefore, is not simply that RSSI is inaccurate, but that prior literature insufficiently theorizes its value within adaptive hybrid pipelines.

\subsubsection{CSI-Based Fine-Grained Localization and Channel Fingerprinting}
In contrast to RSSI, CSI is increasingly treated as the more analytically powerful feature for indoor localization because it preserves sub-carrier-level channel structure and is therefore more responsive to subtle spatial differences. Li et al. (2024) showed that CSI feature fusion can improve localization performance while also supporting anomaly-aware processing, indicating that CSI is valuable not only for precision but also for robustness. Suroso et al.\ (2023) likewise emphasized that CSI-based methods are particularly suited to multipath-heavy environments because they exploit, rather than merely suffer from, rich channel variation \cite{suroso2023csi}. Reyes et al.\ (2025) extended this perspective by arguing that temporal stability is central to CSI fingerprinting performance; rich channel descriptors are only useful if they remain sufficiently consistent for reliable matching over time \cite{reyes2025temporal}.

However, the literature is more persuasive in demonstrating CSI's potential than in resolving its deployment cost. Many studies implicitly assume hardware stability, processing headroom, or calibration routines that are difficult to guarantee on low-cost microcontrollers. This is precisely where the present study positions its contribution. The issue is no longer whether CSI can outperform RSSI in theory; recent literature already suggests that it can. The unresolved issue is whether CSI can be operationalized in a passive, edge-feasible localization framework for hidden devices without inheriting the computational burden of research-grade platforms. Thus, CSI literature supports the study's fine-localization strategy, but it also sharpens the engineering gap that the proposed hybrid adaptive filter must address.

\subsubsection{Hybrid and Low-Cost Edge Localization Architectures}
The tension identified in the preceding literature naturally leads to hybrid and low-cost architectural research. If RSSI is efficient but imprecise, and CSI is precise but computationally heavier, then hybridization becomes less a matter of convenience and more a structural necessity. Mottakin et al.\ (2024), through SWiLoc, showed that ESP32-based Wi-Fi localization can be meaningfully implemented on commodity hardware, thereby challenging the assumption that effective indoor localization must depend on expensive research hardware \cite{mottakin2024swiloc}. Yet this work also underscores a critical limitation: feasibility at the hardware level does not automatically imply suitability for hidden, non-cooperative, multipath-heavy security scenarios. Put differently, low-cost localization has been demonstrated, but low-cost hidden-device localization under severe indoor clutter remains underexplored.

This is where hybrid architectures become analytically important. Li et al. (2024) imply that fusion is not merely an accuracy booster but a mechanism for making richer channel information usable under practical constraints. Even so, the current literature rarely specifies how role separation between RSSI and CSI should be designed in edge-constrained systems. Many hybrid approaches remain descriptive at the feature level but are less explicit about computational scheduling, fallback behavior, and passive operation. The present study addresses this omission by conceptualizing hybridization as an adaptive division of labor: RSSI provides coarse search-area reduction, while CSI is reserved for fine-point localization when higher spatial discrimination is required. This move from feature combination to role-based adaptive processing marks a more explicit response to the implementation gap left by prior work.

\subsection{Related Studies}
Recent empirical studies reveal a clear methodological shift: the field is moving away from proving that indoor localization is possible and toward determining under what constraints it remains reliable. Abdel Ghany et al.\ (2020) are important in this regard because their outdoor LoRa comparison of narrowband CSI and RSSI does more than rank two features; it empirically exposes how environmental noise reshapes the usefulness of each signal representation \cite{ghany2020robustness}. Li et al. (2024) push the literature further by showing that CSI feature fusion can improve localization even under anomalous or imperfect conditions, suggesting that robustness now depends on how features are combined rather than on any single feature alone. Ju et al. (2025) contribute by demonstrating that hidden-device localization is feasible in unknown environments, but their results also imply that the challenge is not simply identification accuracy; it is maintaining localization reliability when the target is passive or concealed.

The same empirical trend is visible in low-cost implementations. Mottakin et al.\ (2024) demonstrate that ESP32-based localization is achievable, while Reyes et al.\ (2025) show that temporal stability remains a major determinant of CSI fingerprinting accuracy \cite{mottakin2024swiloc,reyes2025temporal}. Considered together, these studies suggest that recent progress has been substantial but uneven: one line of work improves signal richness, another improves affordability, and another improves applicability to hidden devices. What remains insufficiently addressed is the integrated case in which all three conditions must be satisfied simultaneously. This is the empirical gap that gives the present study its relevance.

\subsection{Related Systems}
At the systems level, the literature makes clear that design assumptions strongly determine what kind of localization problem can actually be solved. Lumos (Sharma et al., 2022) is influential because it treats hidden-device localization as an operational system challenge rather than a purely algorithmic exercise. However, its broader scope also highlights a limitation for the present study: a system optimized for diverse hidden devices in unfamiliar environments is not necessarily optimized for fine-grained desk-level localization using inexpensive Wi-Fi edge nodes in a fixed educational setting. Thus, Lumos is best understood as a strong systems precedent, but not a direct architectural solution to the present study.

MotionCompass (He et al., 2021) sharpens this contrast further. By relying on movement-direction analysis for hidden wireless camera localization, it demonstrates that specialized systems can outperform more generic positioning pipelines when target behavior is exploitable. Yet this strength is also its limitation from the standpoint of the present research. The current study assumes stationary hidden devices, meaning that motion-derived cues cannot be treated as a reliable source of discrimination. Consequently, MotionCompass helps define the boundary of applicability for behavior-sensitive systems and, by contrast, underscores the need for a signal-centric framework grounded in passive RSSI and CSI measurements.

SWiLoc and related ESP32-based systems are valuable because they establish that low-cost embedded localization is feasible. Nevertheless, feasibility alone is not the same as adequacy. These systems are generally not designed around hidden, non-cooperative targets in severe indoor multipath, nor do they explicitly address adaptive coordination between coarse and fine localization stages. The present study therefore extends the systems literature not by rejecting these precedents, but by combining their affordability lessons with a security-driven localization objective and a more explicit hybrid processing logic.

\subsection{Synthesis of the Review}
The review makes four critical points clear. First, the literature consistently shows that RSSI remains operationally attractive but analytically insufficient for fine-grained indoor localization under severe multipath. Second, CSI is repeatedly validated as the more discriminative signal representation, yet its practical use is constrained by computational overhead, calibration demands, and hardware assumptions that do not always translate well to commodity edge platforms. Third, hidden-device localization research has matured enough to establish the security relevance of the problem, but many system designs still assume conditions that differ from the present study, such as target motion, richer sensing stacks, or broader detection goals that do not prioritize desk-level positioning. Fourth, low-cost localization research has improved accessibility, but it has not fully integrated hidden-device targeting, passive sensing, hybrid feature role separation, and cluttered indoor multipath into a single coherent framework.

These unresolved tensions define the research gap more precisely than a simple claim of ``limited prior work.'' The gap is not that RSSI, CSI, hidden-device systems, or ESP32 localization have been neglected individually. Rather, the gap lies in the absence of a framework that combines them in a way that is both methodologically coherent and practically deployable. Existing studies tend to trade one requirement for another: RSSI-based methods gain efficiency but lose precision; CSI-based methods gain precision but increase computational demand; hidden-device systems improve realism but often depend on assumptions not suitable for static classroom targets; and low-cost systems improve affordability without fully addressing passive, fine-grained localization under heavy multipath. The present study is justified precisely because it responds to this fragmented state of the literature through a hybrid adaptive filter in which RSSI performs coarse localization, CSI performs fine-point discrimination, and ESP32-laptop coordination makes the approach viable within a realistic academic environment. In this way, the study does not merely extend existing literature; it attempts to reconcile several strands of research that remain only partially connected in current scholarship.
